%% Information Processing in Agriculture (Elsevier, ISSN 2214-3173)
%% The journal mandates no specific class; this uses Elsevier's CAS single-column
%% template, which is what the guide-for-authors points at. See templates/.
%%
%% Compile:  latexmk -pdf main.tex
%% Requires: templates/els-cas-templates/{cas-sc.cls,cas-common.sty} on TEXINPUTS,
%%           cas-sc.cls and cas-common.sty are copied beside this file already.

\documentclass[a4paper,fleqn]{cas-sc}

% IPA prints numbered references.
\usepackage[numbers]{natbib}

% Draft scaffolding -- delete both before submission.
\usepackage{xcolor}
\newcommand{\todo}[1]{\textcolor{BrickRed}{\textbf{[TODO:}~#1\textbf{]}}}

\newtheorem{theorem}{Proposition}
\newproof{pf}{Proof}

\begin{document}
\let\WriteBookmarks\relax
\def\floatpagepagefraction{1}
\def\textpagefraction{.001}

\shorttitle{Profit and nitrate together: an RL environment on official SWAT+}
\shortauthors{E. Du et al.}

\title[mode=title]{Profit and nitrate together: a reinforcement-learning environment
on official SWAT+, and a protocol for telling search from adaptation}

\author[1]{Eric Du}[orcid=0000-0000-0000-0000]
\cormark[1]
\ead{ericdu.awesome@gmail.com}
\credit{Conceptualization, Methodology, Software, Formal analysis, Writing -- original draft}

% TODO: fill in -- plain text here on purpose, CAS re-processes these tokens
% and colour macros inside them can break the frontmatter build.
\affiliation[1]{organization={Department},
            addressline={Street},
            city={City},
            postcode={00000},
            state={State},
            country={Country}}

% \todo{co-authors and affiliations to confirm} -- add \author[2]{} blocks here,
% each with its own \credit{} statement (CRediT is required by the journal).

\cortext[1]{Corresponding author}

% Journal limit: 250 words. Currently 225, leaving 25 for the headline finding.
\begin{abstract}
Southern Idaho's irrigated forage systems face two pressures at once: an
over-appropriated aquifer, and dairy manure whose nitrogen is cheap to spread and
costly once it reaches groundwater. Rainfall at Kimberly falls far short of crop
demand, so nearly all water is applied by the grower, and the irrigation calendar is
also the leaching calendar. Machine learning is often proposed for such decisions, but
reported gains are hard to interpret: a learned schedule is usually compared against
routine practice rather than a fixed schedule optimised as hard, and judged on profit
alone, which leaves water and nitrate out of view.

We ask whether a policy that responds to in-season conditions can raise profit without
spending the gain on water and nitrate. Ours is the first reinforcement-learning stack
to wrap the full SWAT+ engine, so water and nitrogen move through the same transport
used in field studies rather than a crop model. Every method receives the same
simulation budget and is evaluated on weather years withheld from training, which
separates a policy that has genuinely adapted from one that was searched harder.
Nitrate is never priced: strategies are compared along a profit--leaching frontier, so
conclusions survive disagreement about what leaching is worth.
% HEADLINE FINDING GOES HERE (25 words max): where each method's frontier lies.
% Source: corrected runs/exp1_*.json
Perfect foresight of the weather is worth $+361$~\$~ha$^{-1}$ at this site (95\,\% interval 220--502 on an effective sample of 1.25 test windows), which bounds what any adaptive policy can recover.
Leaching is unvalidated here and ranks strategies within the model.
\end{abstract}

% Journal limit: 3-5 items, max 85 characters each (spaces included).
\begin{highlights}
\item First Gymnasium environment driven by the official SWAT+ engine on a measured field
\item A protocol separating optimizer search from state-dependent policy adaptation
\item Budgets matched in engine runs, with objective parity tested rather than assumed
\item Nitrate is never priced: the whole profit-leaching frontier is reported instead
\item Weather-specific scheduling is worth 836 \$/ha, so a null result is not vacuous
\end{highlights}

% Journal limit: 1-6 keywords, multi-word keywords discouraged.
% Dropped from the manuscript list to meet the cap: sustainable intensification,
% evaluation protocol, CMA-ES, dairy-forage systems.
\begin{keywords}
SWAT+ \sep reinforcement learning \sep irrigation \sep nitrogen \sep
nitrate leaching \sep multi-objective optimization
\end{keywords}

\maketitle

\section{Introduction}\label{sec:intro}

Southern Idaho's irrigated forage systems face simultaneous pressure on water and
nitrogen. At Kimberly, precipitation is far below potential evapotranspiration, so
nearly all crop water is applied, while dairy manure supplies cheap nitrogen and
elevates nitrate concern in groundwater. The two decisions couple through percolation:
an irrigation schedule is also a leaching schedule. Models without hydrological
transport cannot represent that coupling --- and the one prior SWAT-based RL
environment removes it by assumption, taking ``no/negligible percolation and bypass
flow exiting the soil profile at the bottom'' \citep{madondo2023}. A leaching objective
needs an engine in which leaching can occur.

What managers need is a strategy that performs well --- higher profit where the model
is trustworthy, with resource use and leaching visible --- whether that strategy is a
fixed calendar, a simple feedback rule, or a learned policy. Fixed calendars are easy
to implement; state-dependent rules and reinforcement learning can in principle respond
to mid-season conditions.

Two problems make the existing evidence hard to read. First, \textbf{baselines are
rarely matched}. A learned policy compared against a default schedule, or against a
search given a fraction of the compute, will win for reasons that have nothing to do
with adaptation. Establishing that a policy adapts requires a baseline that has been
searched as hard as the policy was trained, on \emph{the same objective} --- and
objective parity between two code paths is a property to be tested, not assumed.
Second, \textbf{the objective is usually single-valued}. A profit-only reward makes
water and nitrate invisible to the optimizer, so a policy can post a headline gain
while applying more water and leaching more nitrogen than the practice it replaces.
Whether such a policy is an improvement depends entirely on prices the study did not
state.

We therefore treat the \emph{protocol} as a first-class contribution rather than
housekeeping, and report the resource consequences of every strategy alongside its
profit.

\subsection{Contributions}\label{sec:contributions}

\begin{enumerate}
\item \textbf{A Gymnasium environment driven by official SWAT+ rev 62.0.0} on
practitioner \texttt{TxtInOut} inputs and a measured field --- to our knowledge the
first. Strategies act through the same text management files and the same native water
and nitrate routing used in field studies.
\item \textbf{An evaluation protocol that separates search from adaptation}: budgets
matched in engine runs rather than native units, verified objective parity across
methods, year-disjoint weather windows, a perfect-foresight reference, a parametric
feedback controller as a third policy class, and train-selected frozen plans as a
control on state-dependence.
\item \textbf{Profit as the optimized objective, sustainability as a second objective
settled by a frontier rather than a price.} Profit is what the reward maximizes and
what the headline reports. Nitrate and water are not annotations beneath it: each
strategy is re-scored across a swept $\lambda_n$ and the whole profit--leaching
frontier is published, with the profit-only arm always alongside, so no sustainability
conclusion depends on a shadow price we would have to defend. Where a method's frontier
lies wholly inside another's, that is a dominance result; where the frontiers cross, we
report the crossing rather than choosing the $\lambda_n$ that favours our preferred arm.
\item \textbf{A calibration statement that treats per-crop bias, and an unvalidated
percolation pathway, as bounds on interpretation} rather than as caveats appended to a
conclusion.
\end{enumerate}

\section{Related work}\label{sec:related}

\textbf{Agricultural RL.} CropGym, gym-DSSAT, CyclesGym, WOFOSTGym and related
environments wrap process-based crop models for sequential management
\citep{overweg2021, gautron2022, turchetta2022}. Most emphasize crop physiology;
hydrological nitrate transport is often missing or crude. The literature applying RL to
agricultural simulators remains small --- on the order of a dozen papers --- so strong
claims about when closed-loop policies beat fixed schedules are not yet well supported.
Most report a single scalar reward, usually profit or yield, with resource use reported
after the fact if at all.

\textbf{SWATgym and SWAT+.} \citet{madondo2023} introduced SWATgym, an OpenAI Gym
environment for joint fertilizer and irrigation management. Its dynamics are, in the
authors' words, ``the first Python-based implementation of SWAT'' --- a
reimplementation of a subset of classic SWAT \citep{arnold2012}, not the shipped Fortran
executable or a practitioner \texttt{TxtInOut} tree. The paper states the subset
explicitly: inputs ``applied uniformly and daily, over a single growing season
(typically 120 days for corn)''; a ``homogeneous soil profile''; ``no/negligible
percolation and bypass flow exiting the soil profile at the bottom; no lateral and base
flow''; and no nutrient other than nitrogen. Corn is the only crop, and an episode runs
from emergence to harvest with the soil state reset between episodes. SWATgym established
that SWAT-style hydrology--crop coupling is a useful RL testbed, and those simplifications
are what make it fast. They also remove exactly what this study needs: with percolation
assumed away the leaching column is identically zero for every strategy, so a
profit--leaching frontier cannot be drawn; with one season and one crop there is no
manure carry-over, no perennial, and no multi-year soil-nitrogen state to calibrate
against a measured trajectory; and with a hardcoded homogeneous soil there is no input
file into which a site's measured profile, crop coefficients, or management record can be
ported. Our contribution is therefore complementary rather than incremental: wrap
\textbf{official SWAT+} rev 62.0.0 so strategies act through the same text inputs a
hydrologist edits and the engine's own multilayer water and nitrate routing, over a
seven-year rotation replayed from measured practice. The cost is that every evaluation
is a full engine run (\S\ref{sec:protocol}); the return is that a leaching number exists
to report, and that a schedule found here is a SWAT+ management file that can be run
and audited outside the environment. We are not aware of a prior Gymnasium environment
driven by the SWAT+ engine (Table~\ref{tbl:envs}).

\begin{table*}[t]
\caption{Crop-management RL environments (selected).}\label{tbl:envs}
\begin{tabular*}{\tblwidth}{@{}>{\raggedright\arraybackslash}p{0.16\textwidth}>{\raggedright\arraybackslash}p{0.29\textwidth}>{\raggedright\arraybackslash}p{0.09\textwidth}>{\raggedright\arraybackslash}p{0.09\textwidth}>{\raggedright\arraybackslash}p{0.25\textwidth}@{}}
\toprule
Environment & Dynamics engine & Fertilizer / N & Irrigation & Hydrologic N routing \\
\midrule
CropGym \citep{overweg2021} & WOFOST / PCSE & \checkmark & --- & limited \\
gym-DSSAT \citep{gautron2022} & DSSAT & \checkmark & \checkmark & crop-centric \\
CyclesGym \citep{turchetta2022} & Cycles & \checkmark & \checkmark & limited \\
SWATgym \citep{madondo2023} & Python reimplementation of a SWAT subset; one corn season & \checkmark & \checkmark & none: percolation and base flow assumed zero \\
\textbf{This work} & \textbf{Official SWAT+ rev 62.0.0} (\texttt{TxtInOut}) & \checkmark & \checkmark & \textbf{present in engine} (validation bounded, \S\ref{sec:calibration}) \\
\bottomrule
\end{tabular*}
\end{table*}

Table~\ref{tbl:envs} states that routing is \emph{present in the engine}, not that it is
validated at this site; \S\ref{sec:calibration} and \S\ref{sec:limitations} give the
bound.

\textbf{Comparison protocol.} Agricultural optimization and ML papers often report gains
over weak baselines without equalizing search effort or publishing the full candidate
set. Frozen plans --- a policy's train-selected schedule replayed without further
conditioning --- appear occasionally as a diagnostic. We treat them as an object of
study: a frozen plan is only a valid measure of adaptation if it is also a
\emph{well-searched} schedule, and \S\ref{sec:protocol} explains why it usually is not.

\section{Reinforcement-learning formulation}\label{sec:formulation}

Constants below are the values used in the released code.

\subsection{Decision problem}\label{sec:mdp}

Management is a finite-horizon, undiscounted Markov decision process. An episode covers
a seven-year rotation preceded by one spin-up year that \texttt{print.prt}'s
\texttt{nyskip} discards, so a window spans eight calendar years. A step is one
growing-season month, April--September, giving \textbf{42 steps per episode}
(7 years $\times$ 6 months). The discount factor is $\gamma = 1$: the objective is total
rotation profit over a fixed horizon, so discounting would distort it.

\subsubsection{Closing the loop on a batch simulator}\label{sec:prefix-replay}

The MDP above assumes an environment that can be advanced one step at a time. SWAT+
cannot be: it is a batch executable, not a library. It reads a \texttt{TxtInOut}
directory of text inputs, simulates the entire configured period, writes its output
tables, and exits. There is no stepping API, no checkpoint-restart, and no point at
which an action can be injected into a simulation already underway. Management is not
an input \emph{stream} to SWAT+; it is part of the input \emph{file set}, fixed before
the run begins. A naive wrapper therefore gets one action per process --- an eight-year
management plan in, a scalar objective out --- which is an optimization problem, not an
MDP.

We recover the interaction loop by \textbf{prefix replay}. The environment owns the
plan, not the engine. It holds the decision prefix as arrays --- per-month irrigation
depth, per-month mineral N, per-year manure and crop --- initialised to the no-action
state at \texttt{reset}. A step writes the new action into the slot for the current
month, re-renders the \emph{whole} eight-year management from the updated arrays into
\texttt{irrigation.ops}, \texttt{management.sch} and the rest, and re-runs the engine
from the beginning. The observation is then read back out of the monthly output tables
at the row for the month just decided, and the reward is the increment in cumulative
profit that this run produces over the previous one. The agent's step boundary is real;
the simulator's is not.

\noindent\texttt{step($a_t$)} is therefore:

\begin{enumerate}
\item write $a_t$ into the plan at month $t$ (clipped by the remaining-N allowance where
a cap applies);
\item repair rotation feasibility over the whole plan, leaving already-simulated N
applications alone;
\item render plan $\rightarrow$ SWAT+ text inputs, set \texttt{time.sim} to the sampled
weather window;
\item run the engine over all eight years;
\item read \texttt{hru\_wb\_mon} / \texttt{hru\_pw\_mon} at (year, month) $= t$ for the
observation, and the annual tables for profit;
\item return $r_t = (\Pi_t - \Pi_{t-1}) \cdot 10^{-3}$.
\end{enumerate}

\textbf{Why this is the same MDP, and not an approximation of one.} The construction is
only sound if replaying a prefix yields what a pausable engine would have produced. That
rests on two properties of SWAT+, which we state as assumptions because they are
properties of the engine rather than of our code, and which the test suite checks rather
than assumes:

\begin{description}
\item[(A1) Determinism.] Given a fixed input tree and weather window, the engine's
output is a function of its inputs alone --- no sampling, no dependence on wall-clock or
run order.
\item[(A2) Causality in simulated time.] Output rows for month $m$ depend only on inputs
dated at or before $m$. Management operations scheduled after $m$ cannot alter them.
\end{description}

\begin{theorem}\label{prop:prefix}
Let $E(a_0 \ldots a_k)$ denote a run of the engine on the plan whose first $k+1$ monthly
slots are filled and whose remaining slots hold the no-action default, and let
$o_m(\cdot)$ be the observation channels read from the row for month $m$. Under (A1) and
(A2), for every $t$ and every $m \leq t$,
$o_m(E(a_0 \ldots a_t)) = o_m(E(a_0 \ldots a_m))$.
\end{theorem}

\begin{pf}
Fix $m \leq t$. The two runs are given input trees that agree on every operation dated
at or before month $m$; they differ only in slots $m+1 \ldots t$, which carry dates
strictly after $m$. By (A2) the row for month $m$ is invariant to those slots, so both
runs produce the same row for $m$; by (A1) they produce it identically rather than
merely in distribution. \hfill$\square$
\end{pf}

The consequence is the one the protocol needs: the observation the agent received at
step $m$ is never rewritten by anything it does later, so the trajectory
$(o_0, a_0, r_0, o_1, \ldots)$ is exactly the trajectory a checkpoint-restart engine
would have generated. The transition is Markov in the environment's own state --- plan
prefix plus simulated soil state --- even though the engine is stateless between calls,
and the per-step reward telescopes to rotation profit by construction.

Both assumptions are load-bearing enough to test rather than assert. (A2) is checked
directly by seeding the undecided tail of the plan with a large irrigation and requiring
every earlier observation to be bit-identical to the unseeded run
(\texttt{tests/test\_monthly.py::test\_undecided\_tail\_cannot\_change\_earlier\_observations});
the row selection Proposition~\ref{prop:prefix} quantifies over is pinned by
\texttt{test\_observation\_reads\_the\_decided\_month\_not\_the\_last\_row}, and its
consequence --- that the hydrologic channels actually move within an episode --- by
\texttt{test\_hydrologic\_channels\_vary\_within\_an\_episode}. The last of these is not
decoration: a regression that made the observation constant would leave the agent
conditioning on the clock alone, and every adaptation result in
\S\ref{sec:adaptation} would be a foregone conclusion rather than a measurement.

Two properties of this construction matter downstream, and are easy to get wrong:

\begin{itemize}
\item \textbf{The row read must be selected, not taken from the end of the table.}
Because every step simulates all eight years, the monthly tables always end at December
of the final year. Reading the last row returns a \emph{future} month, identical at every
step --- which would both violate the no-foresight constraint below and flatten the
hydrologic channels into constants, leaving the policy nothing to condition on but the
clock. We select by (year, month); the spin-up year that \texttt{nyskip} discards is
accounted for in that index.
\item \textbf{The reward must be an increment, not a level.} Profit is only defined for
a complete rotation, and step $t$'s run reports the whole eight years including the
undecided tail. Taking differences makes the per-step reward the marginal value of the
month just decided, and makes the undiscounted return over an episode equal to the final
rotation profit, which is the quantity \S\ref{sec:reward} prices.
\end{itemize}

The cost of this design is that one step is one full simulation: an episode costs
\textbf{42 engine runs}. This is the single most consequential implementation fact in
the paper --- it is what makes budget matching meaningful across methods that otherwise
count effort in incompatible units, and what makes evaluation expensive enough that
sample efficiency matters.

\subsection{Observation}\label{sec:observation}

The observation is 13-dimensional, \texttt{float32}. All channels describe months
already simulated (Table~\ref{tbl:obs}).

\begin{table}[t]
\caption{Observation channels and scalings.}\label{tbl:obs}
\begin{tabular*}{\tblwidth}{@{}cLL@{}}
\toprule
\# & Channel & Scaling \\
\midrule
0 & Episode progress & $t/42$ \\
1--3 & Current crop, one-hot (corn, barley, alfalfa) & --- \\
4 & Alfalfa stand age & $/7$ \\
5 & Soil water at end of the month just simulated (mm) & $/300$ \\
6 & Nitrogen-stress days, month just simulated & $/50$ \\
7 & Water-stress days, month just simulated & $/50$ \\
8 & Precipitation, month just simulated (mm) & $/50$ \\
9 & PET, month just simulated (mm) & $/200$ \\
10 & Remaining annual N allowance & $/\mathrm{cap}$ \\
11 & Month within season & $/6$ \\
12 & Year within rotation & $/7$ \\
\bottomrule
\end{tabular*}
\end{table}

Divisors are order-of-magnitude scalings chosen to put channels near unit range, not
statistical normalization from a data pass --- and that turned out to be insufficient on
its own, which \S\ref{sec:ppo} treats as a finding rather than an implementation note.
Hydrologic state comes from \texttt{hru\_wb\_mon}, stress from \texttt{hru\_pw\_mon}.

\textbf{No future weather is observable.} The policy is a nowcast, not a forecast. This
is deliberate and it is what makes the perfect-foresight reference in
\S\ref{sec:protocol} a meaningful contrast rather than a trivial one: the gap between
the two quantifies what anticipation would be worth to a manager who cannot anticipate.

\subsection{Action}\label{sec:action}

For the irrigation arm the action is a scalar in $[0,1]$, scaled by
\texttt{MONTH\_DEPTH\_MAX} $= 200$~mm, giving the depth applied in the current month.
Depths become SWAT+ management operations directly: one \texttt{irrm} operation per
non-zero month, placed mid-month at day 15, or on the last day before harvest when
day 15 falls at or after it, written to \texttt{irrigation.ops} with
\texttt{IRR\_EFF} $= 1.0$.

\textbf{Constraints are enforced by repair, not by reward penalty.} An infeasible plan
--- a rotation violation, or nitrogen above a loading cap when one is imposed --- is
rewritten to the nearest feasible plan before simulation. This keeps the reward
interpretable as money, but it has a cost that \S\ref{sec:protocol} makes into a
protocol requirement: repair silently changes what was requested, so a repair applied on
one code path and not another changes the objective invisibly.

Repair is therefore applied \emph{asymmetrically by design}, and the asymmetry is stated
rather than inherited. Rotation feasibility is repaired over the whole plan, since it is
a property of the sequence. A nitrogen cap is \textbf{not}: rewriting the whole plan each
step would let month $t$ revise months already simulated and reported to the agent, so
the cap instead enters as observation channel 10 (remaining annual allowance) and clips
only the current month's application. The agent sees the constraint and acts inside it;
the past is never rewritten. In the irrigation experiment the cap is disabled on every
arm (\texttt{max\_n = None}), and the experiment-level scorers give it \textbf{no default
value} --- it must be passed explicitly --- because a silent default once applied a
400~kg~N~ha$^{-1}$ cap to the open-loop objective while the policy ran uncapped, which
inverted the sign of the headline comparison.

\subsection{Reward}\label{sec:reward}

\textbf{The full reward function, in one place.} The scored quantity is rotation-total
profit, in \$~ha$^{-1}$ over the seven scored years:

\begin{equation}\label{eq:profit}
\Pi(\lambda_n) = \sum_c Y_c \, p_c - \mathrm{mm_{irr}} \, p_w - \mathrm{Mg_{man}} \, p_m
- \mathrm{kg_N} \, p_N - n_\mathrm{events} \, p_\mathrm{op} - \lambda_n \, \mathrm{NO_3}
\end{equation}

Every term is read from a completed engine run rather than assumed
(Table~\ref{tbl:terms}).

\begin{table*}[t]
\caption{Profit terms, their engine sources, and their prices (\S\ref{sec:prices}).}\label{tbl:terms}
\begin{tabular*}{\tblwidth}{@{}LLLL@{}}
\toprule
Term & Symbol & Source & Price \\
\midrule
Crop revenue & $\sum_c Y_c p_c$ & \texttt{basin\_crop\_yld\_yr}, Mg DM by plant name & \$/Mg DM, NASS 2022--24 mean \citep{nass2024} \\
Irrigation & $\mathrm{mm_{irr}} p_w$ & \texttt{hru\_wb\_yr.irr} & \$0.43~mm$^{-1}$~ha$^{-1}$ \citep{wd01_2025,idahopower_sched24} \\
Manure & $\mathrm{Mg_{man}} p_m$ & \texttt{management.sch} fert ops, manure products & \$7.94~Mg$^{-1}$ \citep{isu_custom_2026} \\
Mineral N & $\mathrm{kg_N} p_N$ & same, non-manure products, $\times$ \texttt{fertilizer.frt} N fraction & \$1.54~kg$^{-1}$~N \citep{usda_ams_urea_2025} \\
Application passes & $n_\mathrm{events} p_\mathrm{op}$ & count of fert operations & \$37.07~ha$^{-1}$~event$^{-1}$ \citep{uidaho_custom_2025} \\
Nitrate leaching & $\lambda_n \mathrm{NO_3}$ & \texttt{basin\_aqu\_yr.no3\_rchg} & \textbf{swept}, $\lambda_n = 0$ by default \\
\bottomrule
\end{tabular*}
\end{table*}

Three details matter for reproducing the number. Manure is charged by \textbf{mass}, not
by nitrogen content, because dairy manure here is a disposal stream whose real cost is
haul-and-spread; its N content enters only the emissions accounting
(\S\ref{sec:emissions}). Passes are charged per operation for either source, because
without a per-event charge splitting nitrogen is free and ``how many passes are worth
it'' has no answer. And $\lambda_n$ defaults to \textbf{zero}: leaching is always
\emph{reported} as \texttt{no3\_leached\_kg}, and enters $\Pi$ only when a sweep sets a
price (\S\ref{sec:multiobjective}). N\textsubscript{2}O never enters $\Pi$ at all.

Per step,
\begin{equation}\label{eq:reward}
r_t = (\Pi_t - \Pi_{t-1}) \times 10^{-3}
\end{equation}
where $\Pi_t$ is the profit of the whole rotation re-simulated with operations decided
through step $t$. Two properties follow. The sum telescopes and $\gamma = 1$, so
\textbf{the episode return is exactly $10^{-3} \times$ total rotation profit} --- the
decomposition is a credit-assignment device, not a change of objective, and it is what
lets PPO and CMA-ES be checked for optimizing the same quantity
(\S\ref{sec:protocol}). And a step's reward is not that month's profit but the
whole-rotation change caused by that month's water, which is non-local: April water can
move September yield. This is the direct consequence of SWAT+ having no
checkpoint-restart, and it is why credit assignment here is harder than the 42-step
horizon suggests. The $10^{-3}$ scale exists because unscaled \$~ha$^{-1}$ returns put
the value target near $10^7$.

\subsection{Multi-objective formulation}\label{sec:multiobjective}

\textbf{Profit is the objective the policy maximizes; sustainability is the second
objective, and it is settled by a frontier rather than by a price.} The distinction is
operational, not rhetorical: a sustainability claim that depends on a number we cannot
source is not a claim. Because neither externality has a defensible market price at this
site, both enter as \emph{swept} prices rather than as fixed costs inside the headline
objective:

\begin{equation}\label{eq:frontier}
\Pi(\lambda_w, \lambda_n) = \Pi_0 - \lambda_w \, \mathrm{mm_{irr}} - \lambda_n \, \mathrm{NO_3}
\end{equation}

where $\Pi_0$ is revenue net of the agronomic input costs. The grid is evaluated with
the market-price arm ($\lambda_w$ at the scored water price, $\lambda_n = 0$) always
reported alongside, and we publish the whole curve rather than selecting a point, so no
conclusion rests on a shadow price we would have to defend.

\textbf{Why both axes are swept, and why water is the sharper one.} The farmer's water
price is precisely what \emph{fails} to reflect aquifer scarcity --- which is why
depletion is a problem in the Eastern Snake Plain at all --- so treating it as a fixed
cost would build the sustainability question out of the model. Applied water is
measured, gate-verified against practice to within 0.01\,\%, and strongly responsive to
management; it is the strongest sustainability axis available here.

\textbf{The two axes are not independent, and we say so rather than presenting them as
orthogonal.} In this model leaching is driven by over-irrigation, not by fertilization:
holding irrigation at the measured rate, nitrogen dose produces zero leaching at every
rate tested up to 4{,}490~kg~N~ha$^{-1}$, while raising applied water by 20\,\% produces
32~kg~N~ha$^{-1}$. Both prices therefore act through the same lever, and a
two-dimensional grid shows correlated rather than independent movement.

\textbf{Claims are stated as dominance wherever possible} --- \emph{at matched profit,
strategy X uses Y\,\% less water} --- because a dominance claim requires no price at all
and survives a reader's disagreement about what water or nitrate is worth.

\textbf{Re-scoring is exact; re-optimizing is not.} Profit is linear in every price and
the cost terms are stored separably, so a schedule's profit at any price vector is
arithmetic on stored rows with no further simulation. That answers \emph{how would this
schedule fare if water were dearer}. It does not answer \emph{what schedule would you
choose if it were} --- that requires re-running the search, and the two are reported
separately rather than conflated.

\textbf{What ``leaching'' means here.} NO\textsubscript{3} is \texttt{no3\_rchg} summed
over the scored years --- nitrate reaching the aquifer in recharge. It excludes surface,
lateral and tile pathways. This is the quantity of interest for groundwater loading, and
it is also the quantity most exposed to the percolation limitation in
\S\ref{sec:calibration}.

\subsection{Prices}\label{sec:prices}

\begin{table}[t]
\caption{Prices and their bases. Breakevens are reported wherever a price is
load-bearing; the water price is a composite whose sourced range is stated.}\label{tbl:prices}
\begin{tabular*}{\tblwidth}{@{}>{\raggedright\arraybackslash}p{0.15\linewidth}>{\raggedright\arraybackslash}p{0.2\linewidth}>{\raggedright\arraybackslash}p{0.57\linewidth}@{}}
\toprule
Term & Value & Basis \\
\midrule
Crop prices & Idaho marketing-year averages, 2022--24 & USDA NASS \emph{Crop Values 2024 Summary} \citep{nass2024}; alfalfa:corn ratio 1.489 \\
Alfalfa hay & \$/ton $\div$ ($0.907185 \times 0.88$) & dry-matter conversion \\
Barley & 48~lb/bu, 13.5\,\% moisture & grain-to-DM conversion \\
Corn silage & \$/bu $\times$ $9.0 \div (0.907185 \times 0.35)$ & \textbf{derived} --- no NASS silage series \\
Mineral N & \textbf{\$1.54 / kg N} & USDA AMS urea, Inter-Mountain West, w/e 2025-04-04 \citep{usda_ams_urea_2025}; \$644/short ton at 46\,\% N \\
Water & \textbf{\$0.43 / mm / ha} & rental opportunity cost \$0.235 (Water District 1 common pool, above Milner, \$29.00/acre-ft all-in \citep{wd01_2025}) + sprinkler pumping \$0.20 (Idaho Power Schedule 24, 6.72\,\textcent/kWh + \$16.50/kW \citep{idahopower_sched24}; 45~m head, 0.65 efficiency, 0.6 load factor assumed). Sourced range \$0.235--0.64 (rental only to Tier 7 + pumping); the canal company's own assessment is flat per share and not marginal \\
Manure & \textbf{\$7.94 / Mg} & loading and spreading solid manure, \$7.20/ton, $n = 7$ \citep{isu_custom_2026}; haul-and-spread, not a purchase price \\
Application pass & \textbf{\$37.07 / ha} & dry fertilizer application, \$15.00/acre, $n = 1$ \citep{uidaho_custom_2025}; Iowa \$8.15/acre, $n = 59$ \citep{isu_custom_2026}. Term after CyclesGym \citep{turchetta2022}: without it, splitting N is free \\
\bottomrule
\end{tabular*}
\end{table}

Profit is linear in every price and the cost terms are stored separably, so the price at
which any advantage vanishes is exact arithmetic on stored rows with no further
simulation. We report those breakevens wherever a price is load-bearing,
which converts a result that depends on an unsourced number into a statement of the form
\emph{this holds for any price below X}.

\subsection{Learning algorithm, and why it is configured this way}\label{sec:ppo}

The policy is PPO (\texttt{MlpPolicy}, Stable-Baselines3) over the continuous action
box, trained on four \texttt{SubprocVecEnv} workers --- parallelism is across CPU cores
because SWAT+ is CPU-bound and the engine, not the network, is the cost. Three
configuration choices are not defaults and are load-bearing enough to belong in the
method rather than in a config file.

\textbf{$\gamma = 1$, and rollout lengths keyed to the episode.} The objective is total
rotation profit over a fixed 42-step horizon, so discounting would optimize something
else. \texttt{n\_steps} is 4 episodes (168) and \texttt{batch\_size} 2 episodes (84):
both are set \emph{by formula from the episode length} rather than tuned. This is
deliberate and \S\ref{sec:protocol} relies on it --- hyperparameters that were searched
would need a validation split the weather record cannot support.

\textbf{Observations are statistically normalized, and the statistics travel with the
policy.} With the order-of-magnitude divisors of \S\ref{sec:observation} alone, PPO
learns an \emph{exactly constant} schedule: the standard deviation of applied depth
across held-out weather windows is \textbf{0.000~mm}. Several channels sit near zero in
practice --- nitrogen and water stress divided by 50 when stress runs 0--5 --- and a
network fed them learns to ignore them, which is a failure of input scaling masquerading
as a finding about adaptivity. Running observations through a \texttt{VecNormalize}
filter (clip 10) raises that standard deviation to \textbf{6.4~mm}, raises return by
\textbf{1{,}320~\$~ha$^{-1}$}, and \emph{cuts} applied water from 6{,}116 to 5{,}352~mm.
We report this because a paper whose question is ``does the policy condition on state''
can be answered negatively by this bug alone, and the negative result would look
substantive. The filter's mean and variance are saved beside the policy weights and are
a \textbf{required} argument to every scorer: a policy trained on normalized
observations and scored on raw ones is not a degraded policy but a different objective
--- the same class of silent divergence as the nitrogen-cap mismatch in
\S\ref{sec:action}, and the reason \S\ref{sec:protocol} makes objective parity a tested
property.

\textbf{Initial exploration scale.} \texttt{log\_std\_init} $= -2.0$
($\sigma \approx 0.14$ on the unit action box), against the SB3 default of 0. At
$\sigma \approx 0.37$ the policy applied 6{,}116~mm against a measured 3{,}939 and
saturated 10\,\% of its actions at the box boundary; at $\sigma \approx 0.14$ it applied
4{,}533~mm, earned more, and saturated 2\,\%. A wide initial Gaussian on a bounded action
that is \emph{already} generous --- 200~mm in a single month --- spends the budget in a
region no manager would irrigate in.

Both the observation filter and the exploration scale are hashed into the key that
identifies a saved policy, alongside the price vector, the nitrate price, the cap and
the pinned baseline plan, so a checkpoint cannot silently resume into a run asking a
different question.

\subsection{Emissions accounting}\label{sec:emissions}

Nitrous oxide is computed by IPCC (2019 Refinement) Tier 1 accounting
\citep{ipcc2019} from quantities the simulation already produces --- direct
$\mathrm{EF1} \times N_\mathrm{applied}$, indirect
$\mathrm{EF5} \times \mathrm{NO_3}$ from leaching, and indirect
$\mathrm{EF4} \times \mathrm{FracGASF} \times N_\mathrm{applied}$ from volatilization
--- and reported beside profit. It is \textbf{accounting over simulated nitrogen flows,
not a simulated N\textsubscript{2}O flux}, and is labelled as such throughout.

We do not use the engine's denitrification output as an N\textsubscript{2}O proxy. It is
available and would cost nothing to read, but at this model's calibrated parameters
denitrification is effectively switched off (\texttt{denit\_exp} $= 0.001$ against a SWAT
default near 1.4; \texttt{denit\_frac} $= 1.0$, firing only at saturation), and all
nitrate loss terms together come to under 5~kg~N~ha$^{-1}$ in the worst year. Re-enabling
it would mean re-fitting a nitrogen cycle that is calibrated against a measured soil
nitrate trajectory, trading a validated quantity for an unvalidated one.

In the irrigation experiment nitrogen is pinned, so the direct term is constant across
arms and only the indirect leaching term varies; emissions is therefore reported but not
given its own swept axis, where it would be a rescaled copy of the nitrate axis. It
becomes an independent lever in the nitrogen experiment.

\section{Evaluation protocol}\label{sec:protocol}

Each element exists to prevent a specific way of being wrong.

\textbf{Budgets matched in engine runs.} Timesteps and CMA-ES evaluations are not
comparable units, but both reduce to engine runs: one per environment step, one per
window per open-loop evaluation. Both methods are held to the same total, and the
realized counts are read back off the runner and reported rather than assumed.

\textbf{Objective parity, tested rather than assumed.} Two methods can consume identical
budgets and still optimize different problems if a constraint, price, or repair rule is
applied on one code path and not the other. Because constraints here are enforced by
repair (\S\ref{sec:action}), such a divergence is invisible in the reward --- it changes
what was simulated, not what was priced. We therefore treat parity as a testable
property: an identical plan scored through the open-loop path and through the
environment rollout must agree on profit, applied water, applied nitrogen, yield and
leaching, and this is asserted in the test suite. We recommend the check generally; a
constraint that binds on one arm and not another can invert a headline without producing
any visible anomaly.

\textbf{Year-disjoint weather windows.} Training windows start 1995--2004 and test
windows start 2013--2017; the eight-year spans share no calendar year, and 2012 is left
unused as a buffer. An assertion enforces this at import. The rule is not cosmetic: an
earlier split placed a training window and a test window seven calendar years apart out
of eight, so the optimizer was tuned and evaluated on largely the same weather.

\textbf{Reported uncertainty accounts for window overlap.} Test windows begin in
consecutive years, so adjacent windows share seven of eight calendar years and the
paired differences are strongly correlated. A naive standard error over windows
understates uncertainty. We report intervals from a moving-block bootstrap and state the
effective independent sample size, which for a 12-year test span and 8-year windows is
approximately 1.5 --- considerably smaller than the number of windows suggests. This is
structural: with rotation-length windows and a three-decade weather record, no split
yields many independent test windows.

\textbf{Seed replication with a median representative.} Learned policies are trained at
three or more seeds and the \emph{median} performer is reported. Reporting the best of
$N$ is selection on the evaluation set.

\textbf{Frozen plans selected on train.} A policy is rolled out on each training window,
each rollout yielding a concrete fixed schedule; the schedule with the best training
mean is replayed on the test windows. Selecting the best of $N$ on test would bias the
adaptation estimate negative by construction, as the expected maximum of $N$ noisy
estimates.

\textbf{No validation split, by design rather than by omission.} The training block
carries both roles a validation set would otherwise divide: it is where the optimizers
search, and it is where every selection the protocol makes is resolved --- frozen plans
are chosen on their training means, and the matched-budget open-loop optimum that a
policy is measured against is itself searched on the training windows. We do not carve a
third block out of it. The record cannot support one: eight-year windows over a
1995--2025 weather series yield 10 training and 5 test starts once the buffer year is
removed, and a validation block could only be taken from one of the two, each of which
already sits near an effective independent sample size of 1.5. Selecting a configuration
on two or three overlapping windows is close to selecting at random, so a validation
split of this size would buy the appearance of a safeguard rather than the safeguard.
Nor does the protocol need one: PPO's hyperparameters are fixed by formula from the
episode length rather than searched, seeds are summarized by their median rather than
their best, and CMA-ES and PPO are held to a common budget rather than each tuned to its
own most favourable configuration. A validation set exists to adjudicate a choice among
candidates; we have removed the choices instead. This argument covers policy training
only. It does not extend to calibration, whose parameters were fit against measurements
falling inside both blocks, so the environment itself is not year-held-out from the
evaluation --- a limitation we state in \S\ref{sec:limitations} rather than one this
split resolves.

\textbf{Frozen plans are a state-dependence control, not an adaptation metric.} If a
policy scores no better than its own frozen schedule, it has learned a constant and any
``adaptivity'' is nominal. But the converse does not follow. A policy's emitted schedule
is not a \emph{searched} schedule, so the gap between policy and frozen plan mixes two
effects: how much the policy gains by conditioning on state, and how much worse its
emitted schedule is than a well-optimized one. Where the frozen plan falls below the
matched-budget open-loop optimum, that gap is dominated by search quality and overstates
adaptation. The defensible measure of adaptation is the policy against the
matched-budget open-loop optimum; the frozen plan answers only whether the policy is a
constant.

\textbf{A perfect-foresight reference, reported as an estimate.} Optimizing a separate
open-loop schedule per window and comparing against one shared schedule bounds what
weather-specific information is worth in this action space. It is an \emph{estimate},
not an upper bound: each per-window optimization is itself a finite-budget search over
42 dimensions and can be under-converged, and we have observed a learned policy exceed
it. We report it as \texttt{foresight\_estimate} and never as a ceiling.

\textbf{A parametric feedback controller as a third policy class.} A three-parameter
rule --- $\mathrm{depth} = \mathrm{clip}(a + b \cdot \mathrm{soil\text{-}water\ deficit}
+ c \cdot \mathrm{precipitation\ deficit})$ --- uses the same action space, the same
observations, and the same full-horizon replay, hence the same engine cost per rollout
as the learned policy. It separates two very different conclusions: that adaptation has
little value in this system, or that a particular learning method failed to capture it.

\section{Case study}\label{sec:casestudy}

\subsection{Site}\label{sec:site}

The field is a 1-ha Portneuf silt-loam plot at the USDA-ARS GRACEnet site near Kimberly,
Idaho, under a corn--barley--alfalfa rotation. The scored window is 2013--2019 with 2012
as spin-up. Management comprises 158 measured irrigation events totalling 3{,}938.8~mm
and dairy manure of measured composition, applied in the corn and barley years.
Dry-matter yield is measured in every scored year. A calibrated ArcSWAT model of the same
field supplies secondary references --- ET, percolation, uptake, leaching --- that are
not measured on site and are weighted below measurement.

The generated monthly baseline reproduces measured applied depth to within 0.01\,\%
(3{,}938.9 vs 3{,}938.8~mm), which is asserted before any arm is scored. Without that
gate an irrigation experiment can report gains that are simply a different amount of
water than practice.

\subsection{Calibration: where every number comes from}\label{sec:calibration}

The model is a single 1-ha HRU whose inputs are traced, value by value, to a named source
file; the full audit is in the repository's \texttt{PROVENANCE.md}. What matters for
reading the results is not the fitting procedure --- there was very little of it --- but
which numbers are \emph{measured}, which are \emph{borrowed} from a prior model of the same
field, and which are \emph{fitted}, because each class carries a different kind of
uncertainty. Table~\ref{tbl:sources} sets that out; the rest of this section gives the
numbers a reader needs to interpret the results, and the three structural limits that no
calibration can remove.

\begin{table*}[t]
\caption{Provenance of every model input that affects the reported quantities. MEAS: the
site's primary data (USDA-ARS GRACEnet and Long-Term Manure studies, AgriMet). REF: a
calibrated ArcSWAT/SWAT2012 model of the same field, HRU 119. WB: the collaborator's
crop-parameter workbook. Nothing else in the model has been tuned to an
outcome.}\label{tbl:sources}
\begin{tabular*}{\tblwidth}{@{}>{\raggedright\arraybackslash}p{0.27\textwidth}>{\raggedright\arraybackslash}p{0.3\textwidth}>{\raggedright\arraybackslash}p{0.22\textwidth}>{\raggedright\arraybackslash}p{0.13\textwidth}@{}}
\toprule
Input & Value(s) & Source & Class \\
\midrule
Irrigation events, 2013--2019 & 158 events, 3{,}938.8~mm & GRACEnet irrigation workbooks & measured \\
Manure rate and composition & 4 applications; min/org N and P per source & GRACEnet manure analyses & measured \\
Dry-matter yield (target) & 7 years, all crops & GRACEnet \citep{bierer2022} & measured \\
April soil-nitrate profile (target) & 2013--2019, 0--122~cm, 4 plots & GRACEnet & measured \\
Initial soil NO$_3$ / labile P & 9.33 / 8.13~ppm, mass-weighted April 2013 & GRACEnet soil profile & measured \\
Bulk density & four-plot mean & GRACEnet soil profile & measured \\
Soil profile (AWC, $K_\mathrm{sat}$, C, pH, CaCO$_3$) & Portneuf silt loam, 4 layers to 1{,}220~mm & Long-Term Manure study, \texttt{SWAT Soil Props} & measured \\
Humidity, wind & daily, 2012--2020 & AgriMet TWFI \citep{agrimet} & measured \\
Solar radiation & daily, 2012--2020 & AgriMet via REF & measured \\
Precipitation, temperature & daily, 2012--2019 & AgriMet via RuFaS export & measured \\
Grass-reference ET (target for PET) & AgriMet \texttt{ETOS}, 2013--2020 & AgriMet TWFI \citep{agrimet} & measured \\
Crop coefficients (\texttt{bm\_e}, \texttt{harv\_idx}, \texttt{lai\_pot}, canopy curve, $T_\mathrm{opt}$, $T_\mathrm{base}$, \texttt{rt\_dp\_max}) & alfalfa 17.0 / 1.00 / 5.0; corn silage 50.0 / 1.01 / 4.0; barley 28.0 / 0.52 / 4.0 & WB, \emph{Crop Parameters and Plant Harvest Dates} & collaborator-validated, not fitted \\
Harvest efficiency & 0.98 / 0.54 / 0.95 & WB \texttt{HARV\_EFF} & collaborator-validated \\
Operation dates & planting, harvest, kill, manure 10 April & REF \texttt{.mgt} & borrowed \\
CN2, Manning $n$, slope, slope length & 75, 0.14, 0.012, 122~m & REF \texttt{.mgt}/\texttt{.hru} & borrowed \\
\texttt{esco}, \texttt{epco} & 0.95, 0.50 & REF \texttt{.hru} & borrowed \\
Groundwater, snow, P block & \texttt{alpha\_bf} 0.048, \texttt{melt\_tmp} 2.0, \ldots & REF \texttt{.gw}, \texttt{basins.bsn} & borrowed \\
\texttt{n\_perc}, \texttt{n\_uptake}, \texttt{denit\_exp} & 0.20, 10, 0.001 & REF \texttt{basins.bsn} & borrowed \\
\texttt{pet\_co} & \textbf{0.964} & fitted to AgriMet ETOS, closed form & fitted (1 parameter) \\
\texttt{orgn\_min} & \textbf{0.00091} & fitted to yield + soil-nitrate trajectory & fitted (1 parameter) \\
Weather-generator climatology & station \texttt{IDTWINFALLSWS0} & REF \texttt{.wgn} & borrowed \\
Everything else & 38 files & SWAT+ demo template for the Twin Falls region & default, unexamined \\
\bottomrule
\end{tabular*}
\end{table*}

\textbf{Two fitted parameters, each against one measurement.} Potential ET is linear in
\texttt{pet\_co}, so fitting it to the site's measured grass-reference ET series
(AgriMet TWFI \texttt{ETOS}, 2013--2020) has a closed-form solution: \textbf{0.964}, at which
simulated PET matches measured ET$_0$ to $+0.0$\,\%. The reference ArcSWAT model of the
same field scores $+20.8$\,\% on the same target. \texttt{orgn\_min}, the humus
mineralisation rate, is the only parameter left free after the crop coefficients were
returned to the workbook; it lands at 0.00091 and is weakly identified (the objective
moves by 0.4\,\% across 0.0008--0.0012). A SWAT2012 value of 0.001 borrowed from a site
sweep would have been numerically similar but is not transferable: it reproduces the
measured mineralisation of 209.8~kg~N~ha$^{-1}$~yr$^{-1}$ in SWAT2012 and gives
51.5 here. Every other value is measured, borrowed, or a template default.

\textbf{Crop coefficients are the collaborator's, not ours.} An earlier fit moved nine crop
coefficients and bought 2.3 points of mean absolute yield bias, paid for by making every
alfalfa year 14--17 points worse and by pushing corn's radiation-use efficiency to 64.5
against a site-calibrated 50.0 --- a parameter absorbing a structural artefact rather than
measuring anything. The workbook values were restored and are now asserted by a test.
The price of that decision is visible in Table~\ref{tbl:yield} and is reported rather
than fitted away.

\textbf{Measurement outranks the reference model.} Where the two disagree the measurement
is used: initial soil nitrate (9.33~ppm measured against 0 in the reference), labile P
(8.13 against 5), bulk density, and the 2014 and 2019 irrigation totals. The reference
model's ET, percolation, uptake and leaching are secondary references only, weighted
below measurement and never used as fitting targets.

\textbf{Three structural limits.} \emph{Barley phenology}: typed as a cold annual, an
April-sown crop collapses; it is typed warm annual, which restores a credible heat-unit
trajectory. \emph{Alfalfa residency}: declared perennial, alfalfa remains resident in
years it is not the intended crop, and the engine provides no clean end to that
residency; per-crop agreement within $\pm 15$\,\% is therefore not simultaneously
reachable for all crops. \emph{Legume fixation}: SWAT computes fixation as the soil-supply
shortfall, so soil nitrate and fixation are strict substitutes and the measured pool rise
under alfalfa is unreachable at any parameter setting. The nitrate term is down-weighted
for that reason.

\textbf{Percolation behaves correctly but is not validated here.} Actual ET runs at
0.66--0.70 of PET under measured practice. Percolation is \textbf{threshold-behaved in
applied water}: below crop demand essentially nothing drains, and above it drainage and
nitrate rise together (Table~\ref{tbl:perc}). Implied leachate concentration stays within
8--24~mg~L$^{-1}$ wherever drainage occurs, physically plausible and near the
10~mg~L$^{-1}$ drinking-water standard that motivates regional nitrate concern. The water
and nitrogen pathways are therefore consistent, which is what the frontier in
\S\ref{sec:frontier} requires. But the threshold sits \emph{above} measured practice ---
the measured schedule drains 0--19~mm~yr$^{-1}$ depending on the weather window, and the
generated baseline at the same total drains none --- so $\lambda_n$ can separate
strategies only where one over-irrigates, and a frontier across strategies that all
irrigate sanely passes through zero.

\begin{table}[t]
\caption{Percolation and nitrate response to applied water: the generated monthly
baseline scaled uniformly, 2013 start year, current model. Rule: re-measure with
\texttt{scripts/percolation\_check.py} rather than quote.}\label{tbl:perc}
\begin{tabular*}{\tblwidth}{@{}LRRRR@{}}
\toprule
Scale & Applied (mm, 7 yr) & Percolation (mm yr$^{-1}$) & NO$_3$ (kg N ha$^{-1}$ yr$^{-1}$) & Leachate (mg L$^{-1}$) \\
\midrule
$\times 0.6$ & 2{,}363 & 0.0 & 0.00 & --- \\
$\times 0.8$ & 3{,}151 & 0.0 & 0.00 & --- \\
$\times 1.0$ \emph{(measured depth)} & 3{,}939 & 0.0 & 0.00 & --- \\
$\times 1.2$ & 4{,}727 & 17.4 & 1.45 & 8.3 \\
$\times 1.4$ & 5{,}366 & 77.2 & 9.07 & 11.8 \\
$\times 1.6$ & 5{,}778 & 117.0 & 28.11 & 24.0 \\
\bottomrule
\end{tabular*}
\end{table}

Two limits remain. There is \textbf{no on-site measurement of leaching or percolation},
so reported kg~N~ha$^{-1}$ rank strategies \emph{within the model} and are not
measurements of what this field loses. And in-crop ET runs 15--26\,\% below a regional
OpenET benchmark for the same crops (barley $-15$, corn $-24$, alfalfa $-26$\,\%): the
model has no crop-coefficient mechanism above grass-reference ET, the deficit is partly
water supply rather than canopy, and fitting \texttt{pet\_co} to close it would trade a
measured target for an off-site one. It is reported as a bias, not corrected.

\section{Results}\label{sec:results}

\subsection{Calibration state}\label{sec:calibstate}

Agreement must be read per crop and year (Table~\ref{tbl:yield}). Aggregate
goodness-of-fit is misleading when errors are signed, because they cancel across crops
and across years within a crop.

\begin{table}[t]
\caption{Annual dry-matter yield (Mg ha$^{-1}$), shipped model of 2026-09-10 against the
GRACEnet measurement and the reference ArcSWAT model.}\label{tbl:yield}
\begin{tabular*}{\tblwidth}{@{}LLRRRR@{}}
\toprule
Year & Crop & SWAT+ & GRACEnet & Error & Reference \\
\midrule
2013 & Corn & 18.7 & 22.1 & $-15.4$\,\% & 21.8 \\
2014 & Barley & 9.0 & 5.9 & $+52.4$\,\% & 6.5 \\
2015 & Alfalfa & 13.7 & 9.0 & $+51.9$\,\% & 9.9 \\
2016 & Alfalfa & 21.4 & 14.8 & $+45.0$\,\% & 15.8 \\
2017 & Alfalfa & 24.5 & 16.2 & $+51.4$\,\% & 16.4 \\
2018 & Corn & 20.9 & 22.9 & $-9.0$\,\% & 25.4 \\
2019 & Barley & 6.0 & 8.8 & $-31.2$\,\% & 6.9 \\
\bottomrule
\end{tabular*}
\end{table}

\begin{table}[t]
\caption{Per-crop yield bias against GRACEnet, and the other tier-1 targets.}\label{tbl:bias}
\begin{tabular*}{\tblwidth}{@{}>{\raggedright\arraybackslash}p{0.4\linewidth}>{\raggedright\arraybackslash}p{0.55\linewidth}@{}}
\toprule
Quantity & Value \\
\midrule
Alfalfa yield ($n = 3$) & $\mathbf{+49.2}$\,\% \\
Corn yield ($n = 2$) & $-12.1$\,\% \\
Barley yield ($n = 2$) & $+2.4$\,\% (cancellation: $+52.4$ / $-31.2$) \\
Yield, RMS per-year & 40.4\,\% \\
PET vs AgriMet ET$_0$ & $+0.0$\,\% (1{,}157 vs 1{,}164~mm~yr$^{-1}$) \\
April soil nitrate, RMS miss & 94~kg~N~ha$^{-1}$; 1 of 6 years inside its bracket \\
Humus mineralisation & 51.5 vs 209.8 measured~kg~N~ha$^{-1}$~yr$^{-1}$ \\
\bottomrule
\end{tabular*}
\end{table}

Alfalfa is over-predicted by half in every year, uniformly. That is the cost of the
workbook coefficients (\S\ref{sec:calibration}) and it is carried, not fitted away.
Because alfalfa is three of the seven rotation years and profit is revenue-led, every
dollar \emph{level} in this paper carries that bias. The rotation is fixed and identical
across every strategy compared, so the bias is common-mode in the paired
\emph{differences} the headline claims rest on; it does not cancel in absolute profit or
in anything quoted as \$~ha$^{-1}$ rather than $\Delta$\$~ha$^{-1}$. Barley's crop mean
is cancellation of opposite signs. Corn's shortfall is a growth deficit under the
validated parameters, consistent with its ET running below the regional benchmark.

\subsection{Experiment 1 --- irrigation}\label{sec:exp1}

\textbf{How much is there to adapt to?} Before ranking strategies, the perfect-foresight
reference bounds the question. Optimizing a separate 42-parameter schedule per weather
window beats one shared schedule by $\mathbf{+836 \pm 68}$~\$~ha$^{-1}$ on the held-out
windows ($n = 5$, per-window range $+651$ to $+1{,}062$), far above the
250~\$~ha$^{-1}$ noise floor. Weather-specific scheduling is therefore worth a
substantial amount at this site: a null result for closed-loop control here cannot be
explained by there being nothing to adapt to. Per \S\ref{sec:protocol} this is an
estimate and not an upper bound --- each per-window oracle is a finite-budget search over
42 dimensions and may be under-converged, which biases the figure down.

% Table rows below are being regenerated under verified objective parity (Sec 4).
% The earlier table applied an N loading cap to the open-loop arms (incl. the CMA-ES
% objective) but not to the learned-policy arm, so the two optimized different problems
% at matched budget. Only measured practice, which builds no plan, is unaffected.

\begin{table*}[t]
\caption{Exp 1 test profit (\$ ha$^{-1}$), window starts 2013--2017,
$\lambda_n = 0$.}\label{tbl:exp1}
\begin{tabular*}{\tblwidth}{@{}LRRRRRR@{}}
\toprule
Strategy & Test profit & vs measured & vs open-loop & Water (mm) & NO$_3$ (kg N/ha) & N$_2$O (kg/ha) \\
\midrule
Measured practice & 20{,}925 & 0 & --- & 3{,}939 & --- & --- \\
Generated monthly default & --- & --- & --- & 3{,}939 & --- & --- \\
CMA-ES open-loop & --- & --- & 0 & --- & --- & --- \\
CMA feedback controller & --- & --- & --- & --- & --- & --- \\
PPO policy (median seed) & --- & --- & --- & --- & --- & --- \\
Frozen plan (train-selected) & --- & --- & --- & --- & --- & --- \\
\emph{Foresight estimate (oracle $-$ shared)} & $\mathbf{+836 \pm 68}$ & --- & --- & --- & --- & --- \\
\bottomrule
\end{tabular*}
\end{table*}

% TABLE TO ADD: Exp 1 across PPO seeds.

\subsection{The profit--leaching frontier}\label{sec:frontier}

% RESULTS TO ADD: the lambda_n sweep. Anchor grid points on breakevens computed from
% the lambda_n = 0 results, not a priori: the nitrate price at which the learned policy's
% advantage over open-loop search vanishes, and the price at which its advantage over
% measured practice vanishes. Report both, the full frontier, and the cross-scoring
% matrix (every strategy optimized at one price, re-scored at every other -- exact
% arithmetic, no further simulation).
% NB CLAUDE.md rule 1 (2026-08-28): if the arms do not separate on leaching, that
% non-separation IS the finding -- report it rather than a curve.

The frontier answers the second question, once \S\ref{sec:exp1} has answered the first.
Given a strategy that earns more, does it also reach lower leaching at matched profit
than a schedule searched to the same budget --- that is, does closed-loop control still
pay once the objective becomes timing-sensitive rather than volume-sensitive? A profit
gain bought with more water and more nitrate is not an improvement, and the frontier is
what makes that visible without our having to price either.

\subsection{Adaptation and its baseline}\label{sec:adaptation}

% RESULTS TO ADD: the policy against the matched-budget open-loop optimum (the
% adaptation measure), and against its own train-selected frozen plan (the
% state-dependence control), with the gap between the two interpretations made explicit
% per Sec 4. Bootstrap intervals and effective sample size -- never naive SEs over
% overlapping windows.

\subsection{Limitations}\label{sec:limitations}

A single HRU and soil. Full-horizon replay makes every evaluation cost 42 engine runs,
which bounds achievable budgets. The water price is a composite whose pumping term rests on an assumed head and efficiency, and the application-pass rate rests on a single survey response; breakevens are
reported wherever they are load-bearing. Calibration mixes measured and reference-model
targets without year hold-out, and measured yield, crop identity and manure records end
in 2019, so no independent hold-out period is available. \textbf{Leaching is reported,
and priced in the frontier, but not validated} --- the percolation pathway is a known
rev-62 divergence (\S\ref{sec:calibration}), so a policy optimized against leaching may
be exploiting model error, and that failure would not be visible in the results.
N\textsubscript{2}O is accounting over simulated nitrogen flows, not a simulated flux.
Test windows overlap, so the effective independent sample size is far below the number of
windows. Rebalancing the split to widen the test span, and wiring the site's measured
seasonal water-balance series as an independent check on the model's soil-water storage,
are both available and are left to future work.

\section{Future work}\label{sec:future}

Extend beyond one HRU. Complete the nitrogen, rotation
and joint levers under the same protocol, where emissions becomes an independent axis
rather than a rescaled copy of leaching. Resolve or bound the rev-62 percolation
divergence so that leaching magnitudes can carry field meaning. Validate winning
schedules in an independent simulator.

\section{Conclusion}\label{sec:conclusion}

We set out to ask whether a learned policy can raise farm profit without paying for it in
water and nitrate, and we have built the two things that question needs before it can be
answered honestly: a Gymnasium environment driven by official SWAT+ on a measured
Kimberly field --- to our knowledge the first such stack --- and a protocol that makes a
comparison between a learned policy and a searched schedule mean what it appears to mean.
The protocol's requirements are unglamorous: match budgets in engine runs, verify that
competing methods optimize the same objective rather than assuming it, keep training and
evaluation weather disjoint, report uncertainty that accounts for overlapping windows,
and treat a policy's own frozen schedule as a control on state-dependence rather than as
a measure of adaptation. Without them, ``the policy is greener'' cannot be told apart
from ``the policy was searched harder.''

Profit is what the reward maximizes; sustainability is the second objective, and refusing
to price it is what makes the frontier rather than a scalar the deliverable. We report
where each method's curve lies and where curves cross, never an optimum at a chosen
$\lambda_n$, so a reader who disagrees with us about what nitrate is worth can still read
the result.
% EMPIRICAL FINDING GOES HERE, once the corrected runs land.

Two limits bound how far any of this reaches beyond the model, and both grow rather than
shrink the moment sustainability is treated as a real objective. Leaching ranks
strategies within the model and does not measure what the field loses; optimizing
\emph{for} an unvalidated channel is precisely the regime in which model error is
exploited without becoming visible. And N\textsubscript{2}O is not simulated by SWAT+ at
all --- what we report is IPCC Tier 1 accounting over simulated nitrogen flows, never
priced. A greener policy in these results is a greener policy \emph{in this model}, which
is a claim worth making only because the protocol makes it checkable.

\section*{Declaration of competing interest}

\todo{complete --- required by the journal; state ``The authors declare that they have no
known competing financial interests or personal relationships that could have appeared to
influence the work reported in this paper'' if applicable.}

\section*{Acknowledgements}

\todo{to complete}

\section*{Data and code availability}

Code, model inputs, and experiment scripts: \todo{repository URL}. Input provenance and
build notes: \texttt{PROVENANCE.md} (supplement).

% The journal's research-data policy is Option B: deposit and cite, or state why not.

\printcredits

\bibliographystyle{elsarticle-num}
\bibliography{refs}

\end{document}
